\section{什么是 Generalist Agent：目标、知识与任务跨度}

前一章说明了为何需要主动闭环，本章进一步定义闭环要服务的能力范围。\term{generalist agent（通才智能体）}不是“在一个 benchmark 上覆盖很多关卡”的 specialist，而是能面对不断出现的新目标、调用广泛世界知识，并在大量异质任务之间迁移的系统。这里的\term{open-ended objective（开放式目标）}指任务集合没有固定终点，目标可以被用户、环境或智能体持续提出，并且成功标准未必能由一个手写布尔函数完全描述。

\subsection{从固定 reward 到开放式目标}

经典游戏智能体通常接受一个环境、一个动作空间和一个固定 reward；训练可以非常强，却仍只优化预先定义的目标。开放世界要求系统处理新的对象、新组合、新约束和新的语言描述，因此不能把“玩得很好”直接等同于 generality。课程用三条要求把讨论从宣传词汇变成可检查的设计规范。

\lecturefigure{slide-02-generalist-agent-definition.jpg}{Generalist agent 的三项能力要求}{官方录像恢复幻灯片 V002；00:05:39}

\begin{knowledgebox}{读图：三项要求必须同时成立}
先看左侧 open-ended objectives：任务不能被一张固定清单穷尽；再看中间 world knowledge：系统要知道对象、规则、常识和人类策略；最后看 massively multi-task：能力必须跨大量任务共享。只满足第三项的多任务 benchmark 仍可能是封闭集合，只满足第二项的聊天模型又缺少行动与反馈。
\end{knowledgebox}

若任务来自分布 $\mathcal{T}$，每个任务 $\tau$ 有目标 $g_\tau$ 与回报 $R_\tau$，generalist policy 的训练目标可以抽象为：
\[
\max_\theta\;\mathbb{E}_{\tau\sim\mathcal{T}}\left[\mathbb{E}_{\pi_\theta}\left[\sum_{t=0}^{H_\tau}\gamma^t r_t^{(\tau)}\right]\right],
\qquad \text{并要求对 }\tau'\notin\mathcal{T}_{\text{train}}\text{ 保持迁移能力。}
\]
其中 $\tau$ 表示任务，$g_\tau$ 是语言或多模态目标，$H_\tau$ 是任务时长，$r_t^{(\tau)}$ 是任务相关 reward，$\gamma$ 是时间折扣，$\theta$ 是共享 policy 参数。公式的难点不在求期望，而在 $\mathcal{T}$ 本身会增长、reward 未必可写、测试任务又不能只靠记忆训练清单。

\lecturefigure{slide-03-generalist-agent-requirements.jpg}{构建 generalist agent 的三个 ingredients}{官方录像恢复幻灯片 V003；00:05:45}

\begin{importantbox}{三个 ingredients 是系统依赖关系}
开放环境决定“还能提出什么问题”，internet-scale knowledge base 决定“可以借用哪些人类经验”，foundation model 决定“如何把 observation、语言、记忆与 action 统一表示”。缺少任一环，系统都容易退化为封闭 benchmark、无知识探索或不可迁移的专用控制器。
\end{importantbox}

\teachervoice{讲者先对比 Atari、Go 等 specialist agent，再给出 generalist 的三项要求，是在提醒读者不要用单一高分替代能力定义。一个 agent 即使把固定目标优化到超人水平，只要目标、环境和反馈函数都被冻结，它仍没有回答开放世界中的“下一件值得做的事是什么”。}

\subsection{Ingredient 1：环境复杂度是能力上限}

第一个 ingredient 是\term{open-ended environment（开放式环境）}：对象、地形、规则组合和任务路径足够丰富，智能体不会很快穷尽所有可发现结构。环境不是被动的测试容器；它决定可产生哪些 observation、可执行哪些 action、可构造哪些 curriculum，也决定评估能否暴露长期规划与组合泛化。

\lecturefigure{slide-04-open-ended-environment.jpg}{开放环境决定智能体可以学到的能力上界}{官方录像恢复幻灯片 V004；00:06:18}

\begin{knowledgebox}{读图：环境不是越大越好，而是要有可组合结构}
图中强调 agent capability 被 environment complexity 上界限制。先问环境是否包含对象交互、资源依赖、长期状态变化和多条解法，再问地图面积。纯随机的大空间可能只增加搜索成本；真正有教学价值的是规则与对象能够组合，使旧技能成为新任务的 building block。
\end{knowledgebox}

\lecturefigure{slide-07-minecraft-open-world.jpg}{Minecraft 作为程序生成、可组合的开放世界}{官方录像恢复幻灯片 V007；00:09:39}

Minecraft 的优势不是“流行”本身，而是拥有 procedural generation、采集与制作依赖、战斗、生存、建造和近乎任意的创作目标。读图时先看世界中的资源、工具和空间结构，再看第一人称视角：同一个动作会改变 inventory、地形、危险和后续可达任务。它仍是 simulator，不包含真实机器人全部动力学，却为长期、组合式 agent research 提供了可扩展中间层。

\teachervoice{讲者的表述很直接：agent intelligence 的上限受到 environment complexity 限制。这个判断也提醒我们，若 benchmark 只允许短 horizon、单一 reward 和有限对象，模型再大也可能只是在更快地记住封闭任务，而不是形成开放探索能力。}

\subsection{Ingredient 2：Internet-scale World Knowledge}

第二个 ingredient 是\term{world knowledge（世界知识）}：关于对象属性、规则、配方、因果关系、人类策略和语言约定的可复用先验。对于 Minecraft，它不仅是百科事实，还包括“玩家遇到这个局面通常怎样做”、视频里的操作序列、论坛中的失败经验，以及把自然语言目标拆成技术树步骤的过程知识。

\lecturefigure{slide-05-internet-knowledge-base.jpg}{互联网知识库为智能体提供人类参考手册}{官方录像恢复幻灯片 V005；00:06:57}

\begin{knowledgebox}{读图：knowledge base 不等于网页数量}
数据库图标代表的是可检索、可对齐、可转化的知识。先区分 declarative knowledge（物品是什么）与 procedural knowledge（怎样完成任务），再检查它是否能与 agent 当前 observation、inventory 和 goal 对齐。未经 grounding 的网页文本可能正确，却不能直接变成当前世界里的 action。
\end{knowledgebox}

\subsection{Ingredient 3：Foundation Model 作为统一接口}

第三个 ingredient 是 flexible foundation model。这里的作用不是只生成自然语言，而是把目标、视觉状态、代码、反馈、记忆和 action representation 接进同一个序列建模接口。模型必须既能调用 broad prior，又能接受环境返回的 executable evidence；否则它仍只是离线问答系统。

\lecturefigure{slide-06-foundation-model-for-agents.jpg}{Foundation model 连接开放目标与环境反馈}{官方录像恢复幻灯片 V006；00:07:36}

\begin{warningbox}{“Foundation model”不会自动带来 grounding}
参数规模与 web pretraining 可以提高语言和视觉先验，但当前对象能否被操作、代码能否执行、动作是否安全、reward 是否被钻漏洞，都必须通过 environment interface 和反馈验证。把生成内容直接当作真实能力，是 agent 系统最常见的 category error。
\end{warningbox}

\subsection{本章小结}

Generalist agent 的定义包含开放目标、世界知识与大规模多任务能力；构建它则需要开放环境、互联网知识库和 flexible foundation model。三者分别决定问题空间、先验来源与统一接口，后续 MineDojo、MineCLIP、Voyager、Eureka 和 VIMA 都是在补齐这套 recipe 的不同环节。

